\subsection{Knowledge extraction evaluation}\label{subsec:eval-knowledge-extraction}

The knowledge extraction evaluation scores the pipeline's findings against the silver findings of the calibration cluster (Subsection~\ref{subsec:knowledge-extraction-calibration-set}). These silver findings are produced by sending the same extraction instructions that the voters received, to Opus 4.7, a model more capable than the cascade's best model (Sonnet 4.6). Opus 4.7 is not one of the cascade's voters; therefore, the evaluation is not a replay of the cascade's best tier. The generation process is described in Subsection~\ref{subsec:silver-labels-and-annotation-status}.

\pinktodo{run the sensitivity diagnostic.}

For each case, pipeline findings and silver findings are matched one-to-one. First, six structured fields of each finding -- (claim, subject entity, outcome entity, relation type, direction, category) -- are concatenated, and an embedding value is generated. A pipeline-silver pair is eligible to match only when its embedding cosine similarity is higher than 0.55. Two one-to-one matching strategies are available. The \textbf{greedy} matcher repeatedly takes the unmatched pair with the highest embedding similarity, until no more pairs can be matched. This procedure is simple and deterministic, but not globally optimal, since a greedily selected pair can block a better global pairing. The \textbf{optimal} matcher computes the best global assignment with the Hungarian algorithm, which mitigates the issue that the greedy matcher's selection logic would cause. The matching strategy is a swappable component; the results in Chapter~\ref{chapter:Results} are reported under the optimal (Hungarian) matcher, with the greedy matcher reported as a sensitivity diagnostic.

A matched pair is always counted as a loose true positive; it is also counted as a strict true positive, when none of the three fields -- relation type, direction, category -- disagree between findings. The finding pairs with embedding similarity scores higher than 0.55, which disagree on one or more of these fields, contribute 0.5 to the strict true-positive count. This partial contribution was added, as assigning a score of zero would be too harsh for findings that agree on the underlying claim, but differ in one categorical field.

Loose precision is defined as $n_{\mathrm{matched}} / n_{\mathrm{pipeline}}$, loose recall as $n_{\mathrm{matched}} / n_{\mathrm{silver}}$, and loose F1 as their harmonic mean. Strict precision and recall use the same denominators as their loose counterparts, but use the weighted strict true-positive counts, where relation type, direction, and category fields are also taken into account. Strict F1 is then computed as the harmonic mean of strict precision and strict recall.

Silver labels are LLM-generated and not ground truth; therefore, they inherit the limitations of the silver-generation model. The pipeline configuration is selected on the calibration set of Subsection~\ref{subsec:knowledge-extraction-calibration-set}; the knowledge extraction results reported in Chapter~\ref{chapter:Results} and discussed in Chapter~\ref{chapter:Discussion} are computed on that calibration set, and its generalization to the disjoint held-out evaluation set is reported separately as \textbf{RQ5} (Subsection~\ref{subsec:eval-heldout-generalization}).